% BHU PYQ -- Transcribed & Corrected
% Paper: GGRSE-31 / GGRSE31: Geography of Tourism
% Exam: Under Graduate Semester III Examination, 2025-26 (NEP)
% Category: Skill Enhancement Course (SEC)
% Department: Geography

\documentclass[11pt,a4paper]{article}
\usepackage[a4paper,top=2.5cm,bottom=2.5cm,left=2.8cm,right=2.8cm,headheight=14pt]{geometry}
\usepackage{amsmath,amssymb}
\usepackage[T1]{fontenc}
\usepackage[utf8]{inputenc}
\usepackage{lmodern,microtype,fancyhdr,enumitem,hyperref,lastpage,parskip}

\hypersetup{
    colorlinks=true,
    urlcolor=black,
    linkcolor=black,
    pdftitle={GGRSE31: Geography of Tourism -- BHU Sem III 2025-26 (NEP SEC)}
}

\pagestyle{fancy}
\fancyhf{}
\renewcommand{\headrulewidth}{0.4pt}
\renewcommand{\footrulewidth}{0.4pt}
\fancyhead[L]{\small\textit{Banaras Hindu University}}
\fancyhead[R]{\small\textit{GGRSE-31: Geography of Tourism (SEC)}}
\fancyfoot[C]{\small Page \thepage\ of \pageref{LastPage}}

\newlist{parts}{enumerate}{2}
\setlist[parts,1]{label=(\alph*),leftmargin=*,itemsep=4pt,topsep=2pt}
\newcommand{\pts}[1]{\hfill\textit{[#1]}}

\begin{document}

\begin{center}
  {\large\bfseries BANARAS HINDU UNIVERSITY}\\[2pt]
  {\normalsize Under Graduate --- Semester III Examination, 2025-26 (NEP)}\\[6pt]
  \rule{0.8\linewidth}{0.5pt}\\[6pt]
  {\large\bfseries GEOGRAPHY}\\[4pt]
  {\large\bfseries Paper Code: GGRSE-31 : Geography of Tourism}\\[4pt]
  {\normalsize Time: 2:30 Hours\hfill Maximum Marks: 70}\\[2pt]
  {\small REF NO. CECONF/N/2025-26/089}
\end{center}

\rule{\linewidth}{0.4pt}

\smallskip
\noindent\textit{Instructions: Answer 	extbf{five} questions in all including Question No.~1 which is compulsory. All questions carry equal marks (14 marks each).

\bigskip

\noindent\textbf{Question 1.} Write short notes on the following:\pts{$3.5 \times 4 = 14$}
\begin{parts}
  \item Data sources of tourism.
  \item Tourism as an industry.
  \item Drivers of tourism demand and development.
  \item Concept and principles of sustainable tourism.
\end{parts}

\medskip\rule{\linewidth}{0.2pt}\medskip

\noindent\textbf{Question 2.} Discuss in detail the nature, scope, and key concepts of Geography of Tourism.\pts{14}

\medskip\rule{\linewidth}{0.2pt}\medskip

\noindent\textbf{Question 3.} Describe the salient features of typology and spatial distribution of tourism in the World.\pts{14}

\medskip\rule{\linewidth}{0.2pt}\medskip

\noindent\textbf{Question 4.} Explain in detail the inter-relationship among tourism development, regional planning, and economic growth.\pts{14}

\medskip\rule{\linewidth}{0.2pt}\medskip

\noindent\textbf{Question 5.} Give a comprehensive account of the tourism multiplier effect and tourism forecasting techniques.\pts{14}

\medskip\rule{\linewidth}{0.2pt}\medskip

\noindent\textbf{Question 6.} Elaborate the concept of eco-tourism and its role in environmental conservation with appropriate case studies and examples.\pts{14}

\medskip\rule{\linewidth}{0.2pt}\medskip

\noindent\textbf{Question 7.} Discuss in detail cultural heritage tourism and its conservation strategies in India with appropriate examples.\pts{14}

\vfill
\rule{\linewidth}{0.4pt}
\begin{center}\small
\href{https://bhu-exam.vercel.app/}{Click here to visit our website}\\[4pt]
\href{https://me-aryan.vercel.app/}{Click here to visit our contributor}
\end{center}

\end{document}
